\documentclass[12pt]{article}
\usepackage[utf8]{inputenc}
\usepackage[spanish]{babel}
\usepackage{array}
\renewcommand{\arraybackslash}{\let\\\tabularnewline}

\usepackage{geometry}
\usepackage{graphicx}
\usepackage{fancyhdr}
\usepackage{booktabs}
\usepackage[longtable]{cellspace}
\usepackage{longtable}
\setlength\cellspacetoplimit{5pt}
\setlength\cellspacebottomlimit{5pt}
\geometry{left=1.5cm, right=1.5cm, top=2.5cm, bottom=2.5cm}
\renewcommand{\arraystretch}{1.5}

\begin{document}

\begin{center}
    \textbf{\Large INSTITUTO BÍBLICO DE LAS ASAMBLEAS DE DIOS} \\[0.5em]
    \textbf{\Large PENTATEUCO} \\[0.3em]
    \textbf{CÓDIGO: B1403} \\[0.3em]
    \textbf{MODALIDAD: TUTORÍAS} \\[0.3em]
    \textbf{II CUATRIMESTRE, 2025} \\[1em]
    \textbf{Libro de texto:} \textit{El Pentateuco} de Pablo Hoff y la Biblia \\[1em]
    \textbf{Profesora:} Gloriana Elizondo \hspace{2cm} \textbf{Estudiante:} Gerardo Araya Fallas
\end{center}

\vspace{1cm}

\noindent
\textbf{Tema N\textordmasculine{} 3 \hfill 26 may - 01 jun} \\
\textbf{Título:} Introducción al Éxodo / Egipto Israel viaja a Sinaí \\
\textbf{Temas:} Caída y sus consecuencias, dispersión de las naciones \\
\textbf{Lectura:} Lectura páginas 107 a 138 \textit{El Pentateuco} \\
\textbf{Tarea:} Entrega del diario de lectura


\begin{longtable}{|>{\cellspace\raggedright\arraybackslash}p{5.5cm}|>{\cellspace\raggedright\arraybackslash}p{5.8cm}|>{\cellspace\raggedright\arraybackslash}p{5.8cm}|}
\hline
\textbf{CITA TEXTUAL} & \textbf{¿Qué aprendí? / ¿Por qué me llamó la atención?} & \textbf{Aplicación personal} \\
\hline
«Como Lot escogió la mejor tierra de pastoreo antes que la voluntad de Dios, muy pronto se encontró en Sodoma. Pág. 53» 
&
Es interesante cómo se ve la naturaleza carnal del ser humano reflejada en las decisiones que toma Lot, al escoger lo más visible ante los ojos humanos, sin considerar las consecuencias espirituales.
&
Es importante que Dios esté presente en todas las decisiones que tomamos. Aunque pensemos que por experiencia elegimos bien, solo Dios conoce el futuro y el propósito de cada cosa.
\\
\hline
«Porque su tesoro estaba en Sodoma. Pág. 64» 
&
La esposa de Lot desobedeció la orden de no mirar atrás. Quizás por curiosidad o nostalgia, su desobediencia tuvo consecuencias fatales.
&
Dios nos pide soltar lo que él ha limpiado de nuestras vidas. Mirar atrás es como buscar en la basura aquello que él ya quitó para darnos libertad.
\\
\hline
«El hecho de que Abraham, cuyo cuerpo 'estaba ya como muerto' (Rom. 4:19) haya podido engendrar seis hijos más con Cetura indica que recibió nuevos poderes procreativos al engendrar a Isaac. Pág. 72»
&
Abraham dudó de la promesa de Dios por su edad. A pesar de eso, Dios lo sorprendió con vida nueva, demostrando que nada es imposible para él.
&
Dios no tiene límites. Nuestra fe determina cuánto dejamos que él actúe. No pongamos barreras humanas a su poder.
\\
\hline
«Tenía que aprender las lecciones de fe y consagración. Cada nueva generación necesita aprender por propia experiencia lo que Dios puede hacer por ella. Pág. 76»
&
Dios prueba nuestra relación personal con él. Isaac tuvo que aprender a confiar por sí mismo, no por ser hijo de Abraham.
&
Dios trabaja de forma única con cada uno. Las pruebas nos empujan a buscarle y descubrir su fidelidad de primera mano.
\\
\hline
\end{longtable}

\end{document}
